%%%%%%%%%%%%%%%%%%%%%%%%%%%%%%%%%%%%%%%%%%%%%%%%%%%%%%%%%%%%%%%
%
% LaTeX-pohja LUT-kauppakorkeakoulun kandidaatintutkielmille 
% Versio 1.0. 
% 22.10.2020
%
% Alkuperäisen pohjan tehnyt Tuukka Pohjonen        
% Pohjaa muokannut Maija Hujala                       
%
%%%%%%%%%%%%%%%%%%%%%%%%%%%%%%%%%%%%%%%%%%%%%%%%%%%%%%%%%%%%%%%

\documentclass[a4paper,12pt]{article}
\usepackage[utf8]{inputenc}
% HUOM! Jos kirjoitat kandityösi englanniksi, vaihda alla \usepackage[finnish]{babel} 
% komentoon \usepackage[english]{babel}
\usepackage[finnish]{babel}
%\usepackage[english]{babel}
\usepackage{graphicx}
\usepackage[left=2.5cm,top=3cm,right=2.5cm,bottom=3cm,bindingoffset=0.5cm]{geometry}
\linespread{1.5}
\usepackage{natbib}
\usepackage{multicol}
\usepackage{hyperref}
\usepackage{booktabs}
\usepackage{float}

\begin{document}
\thispagestyle{empty}
\bibliographystyle{agsm}
\setcitestyle{authoryear,open={(},close={)}}

% Kansilehti alkaa tästä
\begin{flushleft}
% Yliopiston logo kuvat-kansiosta
\includegraphics[width=5cm]{kuvat/lut5.jpg}

% HUOM! Jos kirjoitat kandityösi englanniksi, vaihda alta kauppakorkeakoulun nimi yms. englanniksi
\textbf{LUT-kauppakorkeakoulu}\\
Kauppatieteiden kandidaatintutkielma\\
% \textbf{LUT School of Business and Management}\\
% Bachelor's thesis, Business Administration\\

% HUOM! Kirjoita alle kohtaan "Suuntautumisvaihtoehto tähän" oma suuntautumisvaihtoehtosi suomeksi tai englanniksi sen mukaan, kummalla kielellä kirjoitat kandityösi
% Suuntautumisvaihtoehdot suomeksi: Laskentatoimi, Liiketoiminta-analytiikka, Kansainvälinen markkinointi, Strateginen rahoitus, Hankintojen johtaminen
% Suuntautumisvaihtoehdot englanniksi: Accounting, Business Analytics, International Marketing, Strategic Finance, Supply Management
Suuntautumisvaihtoehto tähän
\end{flushleft}

\begin{center}
\vfill
% HUOM! Jos kirjoitat suomeksi, kirjoita ensimmäisten aaltosulkeiden sisään otsikko suomeksi ja toisten englanniksi.Jos kirjoitat englanniksi, vaihda otsikot toisin päin.
\textbf{Kandidaatintutkielman otsikko}\\
\textbf{The title of the Bachelor’s thesis}\\
\vfill
\end{center}

\begin{flushright}
pp.kk.vvvv \\ % tähän palautuspäivämäärä
Tekijä: Etunimi Sukunimi \\ % nimiesi
Ohjaaja: Etunimi Sukunimi \\ % ohjaajasi
% HUOM! Jos kirjoitat kandityösi englanniksi, vaihda alta tekijä ja ohjaaja englanniksi.
% Author: Etunimi Sukunimi \\ % nimiesi
% Supervisor: Etunimi Sukunimi \\ % ohjaajasi

\end{flushright}
\titlepage
\pagenumbering{gobble}

\newpage
% Tiivistelmäsivu alkaa tästä
\noindent
\begin{table}[]
\begin{tabular}{lllll}
\textbf{TIIVISTELMÄ} \\
\\
\\
% HUOM! Vaihda omat tietosi kohtiin: tekijä, tutkielman nimi, suuntautumisvaihtoehto, ohjaaja ja hakusanat.
\textbf{Tekijä}              &  &  &  & Etunimi Sukunimi \\
\textbf{Tutkielman nimi}     &  &  &  & Kandidaatintutkielman otsikko \\
%\textbf{}                    & & & & Otsikon toinen rivi tähän, jos ei mahdu yhdelle \\
\textbf{Akateeminen yksikkö} &  &  &  & LUT-kauppakorkeakoulu \\
\textbf{Koulutusohjelma}     &  &  &  & Kauppatieteet, Suuntautumisvaihtoehto \\
\textbf{Ohjaaja}             &  &  &  & Etunimi Sukunimi \\
\textbf{Hakusanat}           &  &  &  & hakusana1, hakusana2, hakusana3 jne.                              \end{tabular}
\end{table}

\noindent Tähän tulee tiivistelmä

\newpage
% Abstract-sivu alkaa tästä
\noindent
\begin{table}[]
\begin{tabular}{lllll}
\textbf{ABSTRACT} \\                                                                                    
\\                                                                                                  
\\
% HUOM! Vaihda omat tietosi kohtiin: author, title, suuntautumisvaihtoehto englanniksi, supervisor ja keywords
\textbf{Author}           &  &  &  & Etunimi Sukunimi \\
\textbf{Title}            &  &  &  & The title of the Bachelor’s thesis \\
%\textbf{}                    & & & & Otsikon toinen rivi tähän, jos ei mahdu yhdelle\\
\textbf{School}           &  &  &  & LUT School of Business and Management \\
\textbf{Degree programme} &  &  &  & Business Administration, Suunt.vaihtoeht. eng. \\
\textbf{Supervisor}       &  &  &  & Etunimi Sukunimi \\
\textbf{Keywords}         &  &  &  & keyword1, keyword2, keyword3 etc.
\end{tabular}
\end{table}

\noindent Abstract here

\newpage
% Tähän tulee sisällysluettelo automaattisesti
\tableofcontents

\newpage
% Varsinainen kandityö alkaa tästä
\pagenumbering{arabic}
\setcounter{page}{1}
\section{Johdanto}
\subsection{Lähteet, lähdeluettelo ja viittaaminen}
Lähdeluettelo tulee automaattisesti bibliography.bib tiedoston takaa! Eri lähdetietoihin tarvittavat pohjatiedot löytyvät löytyvät googlesta hakusanoilla bibtex ja natbib~\citep{black1973pricing}. 

Toinen vaihtoehto on etsiä tarvittava julkaisu Google Scholarilla. Kun oikea lähde löytyy, klikkaa sen alla olevaa lainausmerkki-ikonia. Avautuu pieni ikkuna, jonka alareunassa lukee "BibTeX EndNote RefMan RefWorks". Klikkaa Bibtex, jolloin aukeaa uusi sivu, jolla ovat julkaisun tiedot. Kopioi kaikki sivulla lukeva ja liitä se bibliography.bib tiedoston loppuun. Kuvallinen ohje jälkimmäiseen proseduuriin löytyy tästä: \href{https://libguides.usask.ca/c.php?g=218034&p=1445751}{BibTex and LaTex From Google Scholar}. 

Suosittelen käyttämään viittauksista kahta eri tyyliä: citep ja citet. Citep tuo viittaukseen tekijöiden nimet sekä vuosiluvun. Citet toimii taas tekstin sisäiseen viittaukseen, esimerkiksi: "\citet{black1973pricing} mukaan liirum laarum". Kirjalähteistä sivunumerot saa näkyviin käyttämällä hakasulkeita sivunumerojen kera kuten tässä esimerkissä: \citep[10-100]{hull2003options}.

\subsection{Kaavat}
Tässä kappaleessa näytetään esimerkki kaavan lisäämisestä. Kaavat tulevat yleensä oman \textit{begin}-komentonsa alle esimerkin omaisesti:
\begin{equation}
f_c = S_0\Phi(d_1) - K \mathrm{e}^{-rT}\Phi(d_2).
\end{equation}
On myös mahdollista sisällyttä kaava tekstin sisään käyttämällä dollarimerkkejä $f_c = S_0\Phi(d_1) - K \mathrm{e}^{-rT}\Phi(d_2)$. Lisää tietoja kaavoista ja niiden kirjoittamisesta löytää esim. Overleafin oppaista.

\subsubsection{Kolmannen tason otsikko}
Lorem ipsum dolor sit amet, consectetur adipiscing elit, sed do eiusmod tempor incididunt ut labore et dolore magna aliqua. Ut enim ad minim veniam, quis nostrud exercitation ullamco laboris nisi ut aliquip ex ea commodo consequat. Duis aute irure dolor in reprehenderit in voluptate velit esse cillum dolore eu fugiat nulla pariatur. Excepteur sint occaecat cupidatat non proident, sunt in culpa qui officia deserunt mollit anim id est laborum.

Lorem ipsum dolor sit amet, consectetur adipiscing elit, sed do eiusmod tempor incididunt ut labore et dolore magna aliqua. Ut enim ad minim veniam, quis nostrud exercitation ullamco laboris nisi ut aliquip ex ea commodo consequat. Duis aute irure dolor in reprehenderit in voluptate velit esse cillum dolore eu fugiat nulla pariatur. Excepteur sint occaecat cupidatat non proident, sunt in culpa qui officia deserunt mollit anim id est laborum."

\subsubsection{Toinen kolmannen tason otsikko}
Lorem ipsum dolor sit amet, consectetur adipiscing elit, sed do eiusmod tempor incididunt ut labore et dolore magna aliqua. Ut enim ad minim veniam, quis nostrud exercitation ullamco laboris nisi ut aliquip ex ea commodo consequat. Duis aute irure dolor in reprehenderit in voluptate velit esse cillum dolore eu fugiat nulla pariatur. Excepteur sint occaecat cupidatat non proident, sunt in culpa qui officia deserunt mollit anim id est laborum.

\section{Kuvat}
Kuvien lisääminen ei ole vaikeaa, mutta kuvien sommittelu paikoilleen voi viedä hieman aikaa. Kuvissa~\ref{fig:figure_levels1} ja~\ref{fig:figure_levels2} on sama mallikuva kahdessa eri koossa. Kuvateksti tulee kuvan alle.
% HUOM! Ohjeita kuvien lisäämiseen esim. täältä: https://www.overleaf.com/learn/latex/Inserting_Images
% HUOM! Ohjeita kuvien paikan säätämiseen esim. täältä: https://www.overleaf.com/learn/latex/Positioning_of_Figures

%  Kuvan lisääminen kuvat-kansiosta. Kuvan leveydeksi valittu tekstin leveys ja paikka määrätty juuri tähän.
\begin{figure}[H]
   \centering
   \includegraphics[width=1\textwidth]{kuvat/mallikuva.pdf}
   \caption{Kirjoita kuvateksti tähän}
   \label{fig:figure_levels1}
\end{figure}

Lorem ipsum dolor sit amet, consectetur adipiscing elit, sed do eiusmod tempor incididunt ut labore et dolore magna aliqua. Ut enim ad minim veniam, quis nostrud exercitation ullamco laboris nisi ut aliquip ex ea commodo consequat. Duis aute irure dolor in reprehenderit in voluptate velit esse cillum dolore eu fugiat nulla pariatur. Excepteur sint occaecat cupidatat non proident, sunt in culpa qui officia deserunt mollit anim id est laborum.

% Kuvan lisääminen kuvat-kansiosta. Kuvan leveys annettu senttimetreinä. Kuvan paikka "jossain tässä lähellä" ja sivun alareunassa.
\begin{figure}[hb]
   \centering
   \includegraphics[width=5cm]{kuvat/mallikuva.pdf}
   \caption{Kirjoita kuvateksti tähän}
   \label{fig:figure_levels2}
\end{figure}

Lorem ipsum dolor sit amet, consectetur adipiscing elit, sed do eiusmod tempor incididunt ut labore et dolore magna aliqua. Ut enim ad minim veniam, quis nostrud exercitation ullamco laboris nisi ut aliquip ex ea commodo consequat. Duis aute irure dolor in reprehenderit in voluptate velit esse cillum dolore eu fugiat nulla pariatur. Excepteur sint occaecat cupidatat non proident, sunt in culpa qui officia deserunt mollit anim id est laborum.

Lorem ipsum dolor sit amet, consectetur adipiscing elit, sed do eiusmod tempor incididunt ut labore et dolore magna aliqua. Ut enim ad minim veniam, quis nostrud exercitation ullamco laboris nisi ut aliquip ex ea commodo consequat. Duis aute irure dolor in reprehenderit in voluptate velit esse cillum dolore eu fugiat nulla pariatur. Excepteur sint occaecat cupidatat non proident, sunt in culpa qui officia deserunt mollit anim id est laborum.Lorem ipsum dolor sit amet, consectetur adipiscing elit, sed do eiusmod tempor incididunt ut labore et dolore magna aliqua. Ut enim ad minim veniam, quis nostrud exercitation ullamco laboris nisi ut aliquip ex ea commodo consequat. Duis aute irure dolor in reprehenderit in voluptate velit esse cillum dolore eu fugiat nulla pariatur. Excepteur sint occaecat cupidatat non proident, sunt in culpa qui officia deserunt mollit anim id est laborum.

\section{Taulukot}
Taulukoiden teko vaatii pientä paneutumista, mutta ei ole mahdotonta. Lorem ipsumeiden alla Taulukossa~\ref{tab:erastaulukko} on malli yksinkertaiselle perustaulukolle. Taulukon otsikko sijoitetaan taulukon päälle.

% HUOM!  Overleafin taulukko-ohjeet löydät täältä: https://www.overleaf.com/learn/latex/Tables
% Tämä työkalu voi myös olla avuksi taulukoiden luomisessa: https://www.tablesgenerator.com/). Lisäksi R, Python ja ehkä myös Stata tukevat paketteja, jotka tulostavat Latex taulukoita. jossakin fiksussa muodossa. 

Lorem ipsum dolor sit amet, consectetur adipiscing elit, sed do eiusmod tempor incididunt ut labore et dolore magna aliqua. Ut enim ad minim veniam, quis nostrud exercitation ullamco laboris nisi ut aliquip ex ea commodo consequat. Duis aute irure dolor in reprehenderit in voluptate velit esse cillum dolore eu fugiat nulla pariatur. Excepteur sint occaecat cupidatat non proident, sunt in culpa qui officia deserunt mollit anim id est laborum.Lorem ipsum dolor sit amet, consectetur adipiscing elit, sed do eiusmod tempor incididunt ut labore et dolore magna aliqua. Ut enim ad minim veniam, quis nostrud exercitation ullamco laboris nisi ut aliquip ex ea commodo consequat. Duis aute irure dolor in reprehenderit in voluptate velit esse cillum dolore eu fugiat nulla pariatur. Excepteur sint occaecat cupidatat non proident, sunt in culpa qui officia deserunt mollit anim id est laborum.

% Taulukko, jonka paikaksi on asetettu h eli "jossain tässä lähellä"
\begin{table}[h]
\centering
\caption{Eräs taulukko}
\begin{tabular}{c c c c} 
\toprule
\textbf{Col1}&\textbf{Col2}&\textbf{Col3}&\textbf{Col4}\\
\midrule
 1 & 6 & 87837 & 787 \\ 
 2 & 7 & 78 & 5415 \\
 3 & 545 & 778 & 7507 \\
 4 & 545 & 18744 & 7560 \\
 5 & 88 & 788 & 6344 \\ [1ex] 
\bottomrule
\label{tab:erastaulukko}
\end{tabular}
\end{table}

Lorem ipsum dolor sit amet, consectetur adipiscing elit, sed do eiusmod tempor incididunt ut labore et dolore magna aliqua. Ut enim ad minim veniam, quis nostrud exercitation ullamco laboris nisi ut aliquip ex ea commodo consequat. Duis aute irure dolor in reprehenderit in voluptate velit esse cillum dolore eu fugiat nulla pariatur. Excepteur sint occaecat cupidatat non proident, sunt in culpa qui officia deserunt mollit anim id est laborum.Lorem ipsum dolor sit amet, consectetur adipiscing elit, sed do eiusmod tempor incididunt ut labore et dolore magna aliqua. Ut enim ad minim veniam, quis nostrud exercitation ullamco laboris nisi ut aliquip ex ea commodo consequat. Duis aute irure dolor in reprehenderit in voluptate velit esse cillum dolore eu fugiat nulla pariatur. Excepteur sint occaecat cupidatat non proident, sunt in culpa qui officia deserunt mollit anim id est laborum.

\newpage
%% Lähdeluettelo tulee automaattisesti tähän
\bibliography{bibliography.bib}

\end{document}